% Lettre de Blaise Pascal à sa sœur Gilberte, 31 janvier 1643 — NAF 15383,
% lecture modernisée du volume entier.
%
% Pass: Opus 5 (claude-opus-5), 17 September 2026 — first pass, unchecked
% against the views by a human, and an interpretation rather than a record.
% Where this file and the transcription differ in substance, the
% transcription (batch-01.fr.tex) is the record.
%
% Sources read: the single batch transcription of this volume and nothing
% else. No image was consulted, and no printed edition of Pascal's
% correspondence — neither the Œuvres of Brunschvicg and Boutroux nor Mesnard
% nor any other — and no other transcription. Where the transcription has a
% gap, this reading says so and does not fill it. Batch 1 covers views 1–20;
% views 21–24 are recorded as skipped in transcripts/status.json (nothing to
% transcribe), so the whole volume is read here.
%
% What the volume turned out to be. Twenty-four views, of which five carry
% writing. Two documents, thirty-one years apart, bound together because one
% person stands in both: Gilberte Pascal. First, a letter of Blaise Pascal
% written from Rouen and dated « Ce Samedy dernier Januier 1643 » (views 13,
% 15, 16), in which a second hand — the father's, Étienne Pascal — has written
% a note of thirteen lines into the blank left in the middle of the second
% page. Second, and not mentioned in the catalogue entry, an oblong leaf
% (views 9–10): a printed-and-filled quittance for the rentes of the Hôtel de
% Ville de Paris, passed at Clermont on 27 June 1674 and signed « G. Pascal »
% by Gilberte as widow and testamentary heir of Florin Périer, with a
% certificate of 4 July 1674 on its verso.
%
% There is no mathematics on these leaves, and this reading does not invent
% any. What it modernises is of four kinds, each footnoted at its first use:
%
%   — the orthography and the abbreviations of 1643 and 1674 (escripre,
%     asseuré, sepmaines, Januier, Moy for « mois », Jl for « il », Vne for
%     « une »; lres for « lettres », ordre for « l'ordinaire », Nores for
%     « notaires », apnt for « à présent », Vre for « votre »);
%   — the fiscal and notarial vocabulary (le voyage des élus, l'ordinaire, la
%     rente constituée, l'état du roi, la ferme générale des aides et entrées,
%     la Cour des aides), given its modern sense in a footnote;
%   — the money of account (livre = 20 sols = 240 deniers) put into decimal
%     livres so that the one computation the two figures of the quittance
%     allow can be done; the division itself is this reading's and the leaf
%     never performs it;
%   — the score notation by twenties of view 10, whose first character the
%     transcription could not read, restored here from the arithmetic alone
%     and said to be so.
%
% Three quantitative statements the leaves themselves make are stated as such,
% with the conditions on which they rest: Blaise's count of the lost letters,
% the father's count of his nights, and the fraction of the rente the crown
% actually funded. The weekday of the date was recomputed.
\documentclass[11pt,a4paper]{article}
% The preamble shared by every transcription of the mathematicians' manuscripts in Gallica.
%
% It rests on one idea: **uncertainty compounds, it does not resolve.** A
% transcription that smooths over a doubtful word has destroyed the only thing
% it was made for — a reader must be able to tell, at every point, what was on
% the page and what was supplied. The apparatus macros below are therefore the
% critical apparatus, not decoration.
%
% The same commands are recognised by `scripts/render.mjs`, which produces the
% site's reading view, and by `scripts/tei.mjs`. A macro added here must be
% added there too, or rendering fails loudly — which is the wanted behaviour.
%
% Comments here are English, like the rest of the repository. The documents
% this preamble sets are French, and that is deliberate: see the skills.

\ProvidesPackage{germain}[2026/09/15 Transcriptions of mathematicians' manuscripts in Gallica]

\usepackage{amsmath,amssymb,amsthm}
\usepackage{mathrsfs}
% Kept from the parent preamble so the renderer's diagram support stays
% available; these manuscripts draw no commutative diagrams, and nothing here uses it.
\usepackage{tikz-cd}
\usepackage[english,french]{babel}
\usepackage{fontspec}

% --- Greek ------------------------------------------------------------------
% The text face stays fontspec's default, Latin Modern, which has no Greek:
% XeTeX drops every glyph it lacks with a line in the log and nothing on the
% page, and the Greek words the manuscripts quote vanished from the PDFs. So
% the Greek and Greek Extended blocks get a XeTeX character class of their own,
% and entering or leaving a run of it switches the family to CMU Serif — the
% same Computer Modern design, with polytonic Greek — and back. Only the
% family changes: a Greek word in \textit or \textbf keeps its shape and series.
% The transitions to and from every other class carry the tokens that class
% has towards an ordinary letter, so French spacing before « ; » or inside
% « » still holds after a Greek word. No grouping, so no brace or \endgroup
% can come out unbalanced; maths is left alone, since XeTeX inserts nothing
% there. Kept identical in `exercices.sty`.
\makeatletter
\newfontfamily\germainGreekfont{cmun}[
  Extension=.otf, NFSSFamily=germaingreek,
  UprightFont=*rm, ItalicFont=*ti, SlantedFont=*sl,
  BoldFont=*bx, BoldItalicFont=*bi, BoldSlantedFont=*bl]
\newXeTeXintercharclass\germain@greek
\def\germain@greekrange#1#2{%
  \@tempcnta=#1\relax
  \loop
    \XeTeXcharclass\@tempcnta=\germain@greek
    \ifnum\@tempcnta<#2\advance\@tempcnta\@ne\repeat}
\germain@greekrange{"0370}{"03FF}
\germain@greekrange{"1F00}{"1FFF}
\def\germain@greekprev{\rmdefault}
\def\germain@greekin{%
  \edef\germain@greekprev{\f@family}\fontfamily{germaingreek}\selectfont}
\def\germain@greekout{\fontfamily{\germain@greekprev}\selectfont}
\def\germain@greekjoin{%
  \XeTeXinterchartokenstate=\@ne
  \@tempcnta=\z@
  \loop
    \ifnum\@tempcnta=\germain@greek\else
      \XeTeXinterchartoks\germain@greek\@tempcnta=\expandafter{%
        \expandafter\germain@greekout\the\XeTeXinterchartoks\z@\@tempcnta}%
      \XeTeXinterchartoks\@tempcnta\germain@greek=\expandafter{%
        \the\XeTeXinterchartoks\@tempcnta\z@\germain@greekin}%
    \fi
    \ifnum\@tempcnta<4095\advance\@tempcnta\@ne\repeat}
% babel-french sets its own classes, and saves and restores the interchar
% state, each time a language is selected: join again after it does.
\addto\extrasfrench{\germain@greekjoin}
\addto\extrasenglish{\germain@greekjoin}
\AtBeginDocument{\germain@greekjoin}
\makeatother

\usepackage{geometry}
\geometry{a4paper,margin=25mm,marginparwidth=28mm}
\usepackage{marginnote}
\usepackage{xcolor}
\usepackage[normalem]{ulem}
\setlength{\emergencystretch}{3em}
\usepackage{hyperref}
\hypersetup{colorlinks=true,linkcolor=black,urlcolor=black,pdfborder={0 0 0}}

% --- Batch metadata ---------------------------------------------------------
% Taken as they stand by the rendering script, which shows them at the head of
% the reading view. They fix what the file claims to cover. The macro names are
% the parent project's, so that the scripts stay shared; \folder here means the
% volume's slug — `fr-9115` — and \pages counts Gallica views.

\newcommand{\folder}[1]{\gdef\grothFolder{#1}}
\newcommand{\batch}[1]{\gdef\grothBatch{#1}}
\newcommand{\pages}[2]{\gdef\grothFirst{#1}\gdef\grothLast{#2}}
\newcommand{\foldertitle}[1]{\gdef\grothFoldertitle{#1}}

% The holder's own shelfmark, printed as they write it: « Français 9115 ».
\newcommand{\shelfmark}[1]{\gdef\grothShelfmark{#1}}

% The dating, copied verbatim from the holder's catalogue — never inferred
% here. « XIXe siècle » is the BnF's; a pass that dates a leaf from its content
% says so in a \note{}, as its own inference.
\newcommand{\dating}[1]{\gdef\grothDating{#1}}

% The status notice, printed in the title block and repeated as a diagonal
% watermark on every page of the PDF. The manuscripts are in the public
% domain; what this line declares is not a right but a quality — an
% unverified first pass by a machine. The skills write the line; a file
% without it gets no watermark, so leaving it out is visible in the output.
\newcommand{\watermark}[1]{\gdef\grothWatermark{#1}}

\gdef\grothFolder{?}\gdef\grothBatch{}\gdef\grothFirst{?}\gdef\grothLast{?}%
\gdef\grothFoldertitle{}\gdef\grothShelfmark{}\gdef\grothDating{}\gdef\grothWatermark{}

% --- The résumé -------------------------------------------------------------
\newenvironment{resume}
  {\small\setlength{\parskip}{0.45em}}
  {\par\medskip\hrule\medskip}

% --- Keywords ---------------------------------------------------------------
% In English, closing the modernised reading's résumé: the modern vocabulary
% under which to search for what the leaves construct. The manifest extracts
% them to tag the volume — this line is the single source of the tags.
\newcommand{\keywords}[1]{\par\smallskip\noindent
  {\footnotesize\textsc{Keywords} --- \textit{#1}}\par}

% --- The critical apparatus -------------------------------------------------

% The Gallica view number, in the margin. This is the anchor that lets a
% disagreement over a formula be settled by looking at *one* image rather than
% a volume of 750. The site uses it to turn the facsimile as the transcription
% is scrolled. A view is one scanned image: usually one side of a leaf.
\newcommand{\page}[1]{%
  \marginnote{\footnotesize\color{gray!70}\textsf{v.\,#1}}%
  \label{page:#1}\ignorespaces}

% The BnF's pencilled foliation, where the leaf carries it and a pass can read
% it: « 198r ». This is what the literature cites — « ff. 198r–208v » — and
% the one line that turns a citation into a view. Recorded, never inferred:
% a folio a pass cannot read is not written.
\newcommand{\folio}[1]{%
  \marginnote{\footnotesize\color{gray!50}\textsf{f.\,#1}}\ignorespaces}

% The modernised reading's coarser counterpart to \page{}: which views a
% section is drawn from.
\newcommand{\pagerange}[2]{%
  \marginnote{\footnotesize\color{gray!70}\textsf{v.\,#1--#2}}%
  \ignorespaces}

% Illegible: never guessed.
\newcommand{\ill}{\textcolor{red!60!black}{[\,\ldots\,]}}

% A doubtful reading: the word is given, and flagged as uncertain.
\newcommand{\uncertain}[1]{\uline{#1}}

% An editorial addition — a missing letter, an implied word.
\newcommand{\add}[1]{\textcolor{blue!50!black}{[#1]}}

% Struck out by the writer. What was crossed out is part of the document.
\newcommand{\struck}[1]{\sout{#1}}

% The transcriber's note, on the state of the page or the mathematics.
\newcommand{\note}[1]{\par\smallskip{\small\color{gray!60!black}\textit{[#1]}}\par\smallskip}

% A marginal note of hers, distinct from the body of the text.
\newcommand{\marginal}[1]{\par\smallskip\noindent\hspace{1em}%
  {\small\color{gray!60!black}#1}\par\smallskip}

% --- Title ------------------------------------------------------------------
% The title block is French, like the documents themselves.

\AddToHook{shipout/background}{%
  \ifx\grothWatermark\empty\else
    \begin{tikzpicture}[remember picture, overlay]
      \node[rotate=52, scale=3.4, align=center, text=gray!18,
            font=\sffamily\bfseries]
        at (current page.center) {\grothWatermark};
    \end{tikzpicture}%
  \fi}

\AtBeginDocument{%
  \begin{center}
    {\large\bfseries Manuscrits de mathématiciens (Gallica) ---
      \ifx\grothShelfmark\empty\grothFolder\else\grothShelfmark\fi}\\[2pt]
    {\itshape\grothFoldertitle}\\[4pt]
    {\small vues \grothFirst--\grothLast
      \ifx\grothBatch\empty\else\ (lot \grothBatch)\fi}\\[2pt]
    \ifx\grothDating\empty\else
      {\small\color{gray!60!black}Datation du catalogue : \grothDating}\\[2pt]
    \fi
    \ifx\grothWatermark\empty\else
      {\small\bfseries\color{red!45!gray}\grothWatermark}\\[2pt]
    \fi
    {\footnotesize\color{gray!60!black}Transcription établie d'après les images IIIF de Gallica.
     Source gallica.bnf.fr / Bibliothèque nationale de France.\\
     Les lectures incertaines sont soulignées, les illisibles notées [\,\ldots\,],
     les ajouts entre crochets ; \textsf{v.} numérote les vues de Gallica,
     \textsf{f.} la foliotation au crayon de la BnF.}
  \end{center}
  \vspace{1em}}

\endinput


\folder{naf-15383}
\shelfmark{NAF 15383}
\pages{1}{20}
\dating{}
\watermark{Lecture automatique\\interprétation non vérifiée}
\foldertitle{Lettre de Blaise Pascal à sa sœur Gilberte, 31 janvier 1643 — lecture modernisée du volume entier}

\begin{document}

\section*{Résumé}

\begin{resume}
Ce volume ne contient pas de mathématiques. C'est une lettre de famille de
vingt-cinq lignes, écrite de Rouen en janvier 1643 par un garçon de vingt ans
à sa sœur aînée, et à laquelle leur père a ajouté un billet dans le blanc
resté au milieu de la page ; on y a joint, bien plus tard et sans que le
catalogue le dise, une pièce comptable de 1674 que la même sœur a signée. Le
garçon est Blaise Pascal, le père Étienne Pascal, la sœur Gilberte, mariée à
Florin Périer et appelée dans l'adresse « Mademoiselle Périer la
conseillère ». Rien dans ces feuillets n'annonce l'œuvre : ni la machine
arithmétique, que Pascal construisait alors précisément pour les comptes que
la lettre évoque, ni la géométrie, ni rien de ce qui viendra. Ce qu'ils
donnent est l'envers du travail — la tournée des commissaires, le courrier qui
se perd, une sœur malade dont on n'a pas de nouvelles, un père qui ne se
couche pas.

Le problème de la lettre, s'il faut en nommer un, est un problème de courrier.
Gilberte s'est plainte d'un reproche ; elle affirme écrire à Rouen une fois par
semaine. Blaise répond par un compte : il n'en reçoit pas une toutes les trois
semaines. Il en tire, sous condition, la seule conclusion possible — les
lettres se perdent — et cette conclusion, si les deux affirmations sont
exactes, chiffre la perte : plus de deux lettres sur trois n'arrivent pas.
Le père, de son côté, mesure sa charge de la même manière, par un dénombrement
de nuits : en quatre mois il ne s'est pas couché six fois avant deux heures du
matin, soit moins d'une nuit sur vingt. C'est la façon dont cette famille dit
l'excès : en comptant.

La pièce de 1674 est un autre genre de compte, et elle en cache un. C'est une
quittance de rente sur l'Hôtel de Ville de Paris : Gilberte, devenue veuve,
reconnaît avoir reçu cent trente-trois livres quatre sols pour le second
semestre de l'année, à cause d'une rente de huit cent quatre-vingt-huit livres
constituée en 1635. La pièce donne les deux nombres et ne les rapproche
jamais. Si on les rapproche, le semestre dû est de quatre cent quarante-quatre
livres et le semestre payé de cent trente-trois livres quatre sols, exactement
trois dixièmes du dû. Le roi ne servait plus que trente pour cent de la rente ;
et le compte tombe juste au sol près parce que le cinquième d'une livre fait
quatre sols dans une monnaie de compte à vingt sols. La même vingtaine
organise, sur le verso, le chiffre écrit à l'ancienne — celui dont on ne lit
plus que la fin.

Où cela mène : non pas aux mathématiques de Pascal, mais à leur occasion et à
leur entour. Du côté de l'homme, aux années de Rouen, à la charge de
commissaire pour la levée des tailles en Haute-Normandie, et à Gilberte Périer,
qui écrira la vie de son frère. Du côté des choses, à l'histoire
de la poste d'Ancien Régime et de son \emph{ordinaire}, à celle des rentes
constituées sur l'Hôtel de Ville et des retranchements que la monarchie leur
faisait subir, à la monnaie de compte en base vingt et douze, et à la notation
par vingtaines, qui a laissé « quatre-vingts » dans la langue.

\keywords{Blaise Pascal, Gilberte Périer, Étienne Pascal, Rouen 1643, family
correspondence, seventeenth-century French orthography, notarial act, rente
constituée, Hôtel de Ville de Paris annuities, livre tournois, vigesimal
numerals, money of account}
\end{resume}

\section*{Ce que le volume contient}

\pagerange{1}{20}
Vingt-quatre vues, dont les vingt premières ont été transcrites et les quatre
dernières écartées comme ne portant rien.\footnote{Le fichier
\emph{transcripts/status.json} porte « naf-15383\#2 » à \emph{skipped}. Les
vingt vues transcrites se répartissent ainsi : vues 1 à 8, la reliure, la
chemise marbrée et trois gardes blanches ; vues 9 et 10, un feuillet à
l'italienne ; vues 11 et 12, des feuillets blancs ; vue 13, la première page
de la lettre ; vue 14, son verso, blanc, où le recto transparaît ; vue 15, la
seconde page écrite, avec la signature et le billet de la seconde main ; vue
16, le panneau d'adresse ; vues 17 et 18, des feuillets blancs ; vues 19 et
20, la chemise dépliée et le plat de derrière.} Cinq vues portent de
l'écriture. Aucune foliotation de la bibliothèque n'est lisible nulle part,
de sorte que ni la transcription ni cette lecture ne citent de feuillet :
les renvois se font aux vues de Gallica.\footnote{Ce que les feuillets portent
en fait de chiffres modernes est la cote au crayon « N. a. fr. 15383 »,
soulignée, en haut de la vue 13 et au bas de la vue 9 ; le cachet ovale
« ACHAT N\textsuperscript{o} 24199 » sur les deux ; et, dans la marge gauche de
la vue 9, trois marques de greffe à l'encre dont le chiffre 84. La notice ne
donne pas de datation pour le volume, et le champ correspondant est laissé
vide ici comme dans la transcription, bien que les deux pièces se datent
elles-mêmes.}

Le volume réunit deux documents que rien ne relie qu'une personne. Le premier
est une lettre de janvier 1643 ; le second une quittance de juin 1674, que le
catalogue ne mentionne pas. L'adresse de 1643 porte « Mademoiselle Périer la
conseillère » ; l'acte de 1674 est passé par « dame Gilberte Pascal, veuve et
héritière testamentaire de défunt M\textsuperscript{e} Florin Périer, vivant
conseiller du roi en sa Cour des aides de Clermont-Ferrand », et signé de sa
main. C'est la même femme, trente et un ans plus tard, désignée les deux fois
par la charge de son mari — d'abord comme sa femme, puis comme sa
veuve.\footnote{L'identification tient aux seuls feuillets : le nom, la ville
de Clermont, la charge de conseiller. La lettre ne nomme jamais sa
destinataire autrement que « ma chère sœur » ; c'est l'adresse de la vue 16
qui donne le nom.}

\section*{La lettre de Rouen (vues 13, 15 et 16)}

\pagerange{13}{16}
La lettre s'ouvre en haut de la vue 13 par la date et l'appel, court jusqu'au
bas de la page où elle s'interrompt en pleine phrase, et reprend en haut de la
vue 15 pour s'achever sur la signature. Le verso de la vue 13 est resté blanc.
Le feuillet a été plié et l'adresse écrite au dos du second (vue 16) : c'est
la lettre elle-même qui fait son enveloppe, comme il est d'usage avant les
enveloppes.

\subsection*{Le texte en français d'aujourd'hui}

De Rouen, ce samedi dernier janvier 1643.\footnote{Le feuillet porte « De
Rouen Ce Samedy dernier Januier 1643 », sur deux lignes en tête à droite ;
les deux derniers chiffres du millésime sont tracés en grand et descendent
sous la ligne. « Dernier janvier » se lit ici comme le dernier jour du mois.}

Ma chère sœur,

Je ne doute pas que vous n'ayez été bien en peine du long temps qu'il y a que
vous n'avez reçu de nouvelles de ces quartiers-ci. Mais je crois que vous vous
serez bien doutée que le voyage des élus en a été la cause ; comme en effet,
sans cela, je n'aurais pas manqué de vous écrire plus souvent.\footnote{Le
feuillet a « le voyage des Esleuës ». Le mot désigne, dans le vocabulaire
fiscal du temps, les officiers d'une \emph{élection}, circonscription de la
levée des tailles ; « le voyage des élus » est la tournée que les commissaires
du roi faisaient dans ces circonscriptions. La graphie du feuillet ne permet
pas de trancher entre les officiers et les circonscriptions, et la traduction
proposée ici retient les officiers parce que la phrase suivante nomme les
commissaires.} J'ai à te dire que, messieurs les commissaires étant à Gisors,
mon père me fit aller faire un tour à Paris, où je trouvai une lettre que tu
m'écrivais, où tu me mandais que tu t'étonnais de ce que je te reproche que tu
n'écrivais pas assez souvent ; où tu me dis que tu écris à Rouen toutes les
semaines une fois.\footnote{« Gizors » sur le feuillet : Gisors, dans le Vexin
normand. Le père de l'écrivain est ici l'un de ces commissaires, ou du moins
se déplace avec eux ; le feuillet ne dit que « mon Pere ». Étienne Pascal
était alors commissaire député par le roi pour l'impôt et la levée des tailles
en Haute-Normandie — ce fait n'est pas sur le feuillet et vient de l'histoire
générale, non d'une édition de cette lettre.}

Il est bien assuré, si cela est, que les lettres se
perdent,\footnote{L'abréviation du feuillet est « l\textsuperscript{res} », le
groupe \emph{res} surmonté d'un trait : la forme ordinaire de « lettres ».}
car je n'en reçois pas toutes les trois semaines une.\footnote{Le verbe est
écrit vite et la transcription donne « reçois » avec doute, la cédille n'étant
pas sûre.}

Étant retournés à Rouen, j'y ai trouvé une lettre de M. Périer, qui mande que
tu es très malade. Il ne mande point si ton mal est dangereux, ni si tu te
portes mieux, et il s'est passé un ordinaire depuis sans avoir rien de
lettres ;\footnote{Deux choses ici. Le mot rendu par « très » est, sur le
feuillet, bref, écrit en partie au-dessus de la ligne, surmonté d'un trait
d'abréviation et posé sur une pliure : la transcription l'offre avec doute, et
la lecture reste incertaine. « Ord\textsuperscript{re} » est l'abréviation
usuelle de « l'ordinaire », c'est-à-dire le courrier régulier ; un ordinaire
passé est une levée manquée.} tellement que nous sommes en une peine dont je
te prie de nous tirer au plus tôt.\footnote{La phrase s'interrompt au bas de
la vue 13 sur « Tellement que Nous Sommes » et reprend en haut de la vue 15
sur « En Vne peine ». Le bas du feuillet de la vue 13 est blanc et taché.}
Mais je crois que la prière que je te fais ici sera inutile, car avant que tu
aies reçu cette lettre-ci, j'espère que nous aurons reçu lettres ou de toi ou
de monsieur Périer.\footnote{Le feuillet a « la prie », suivi d'un long trait
horizontal, que la transcription interprète comme une abréviation possible de
« priere ». Avant « Nous aurons », deux mots biffés, « ou tu ».} Le départ,
et le fâcheux du mois.\footnote{Le feuillet porte « Le depart », puis un mot
de cinq ou six jambages qui traverse une pliure et n'a pas été lu, puis « du
fascheux du Moy », net. « Moy » est ici le mois. La phrase reste incomplète :
ce que cette lecture en donne n'est pas une restitution mais l'ordre des mots
lus, et le sens n'en est pas établi.}

Si je savais quelque chose de nouveau, je te le ferais savoir. Je suis,

Ma chère sœur, votre très humble et très affectionné serviteur et frère,

Pascal.\footnote{« Ma Chere Seur » est écrit en grand à gauche au bas de la
page, la formule et la signature en bas à droite. La signature est grande et
suivie d'un paraphe. Le feuillet ne porte que ce nom ; que le signataire soit
Blaise tient à ce qu'il se dit le frère de la destinataire et parle de « mon
Père » à la troisième personne.}

\subsection*{La date, et ce qu'elle vérifie}

« Ce samedi dernier janvier 1643 » se laisse vérifier : le 31 janvier 1643
était un samedi, et c'était aussi le dernier samedi du mois.\footnote{Calcul
fait pour le calendrier grégorien, en usage en France depuis 1582. Les samedis
de janvier 1643 tombent les 3, 10, 17, 24 et 31. Les deux lectures possibles
de la formule — « le dernier jour de janvier, qui est un samedi » et « ce
samedi, dernier samedi de janvier » — désignent donc le même jour, ce qui
n'arrive que lorsque le mois finit un samedi. La date de la notice est ainsi
confirmée par le feuillet seul.}

\subsection*{Le compte des lettres perdues}

La lettre contient un raisonnement, et c'est le seul du volume. Gilberte
affirme écrire à Rouen une fois par semaine ; Blaise constate qu'il n'en reçoit
pas une toutes les trois semaines. Il pose la condition — « si cela est » — et
conclut que les lettres se perdent. La conclusion est juste, et elle se
chiffre : si l'on écrit une lettre par semaine et qu'il en arrive moins d'une
par trois semaines, c'est que plus de deux envois sur trois n'arrivent
pas.\footnote{Le feuillet ne fait pas cette division et ne donne aucun taux :
le rapport de deux tiers est celui de cette lecture, et il ne vaut que sous
les deux affirmations telles qu'elles sont rapportées — celle de Gilberte, que
Blaise ne fait que répéter d'après sa lettre, et la sienne. Rien n'oblige
d'ailleurs à croire l'une ou l'autre exacte : c'est une querelle de frère et
sœur, et le compte est un argument dans la querelle.} La lettre est donc, en
son fond, l'établissement d'une perte par comparaison de deux fréquences. Le
reste — la maladie, l'ordinaire manqué, la peine — vient de ce que le même
défaut de courrier empêche d'avoir des nouvelles quand elles pressent.

\subsection*{Le passage du « vous » au « tu »}

La lettre commence au \emph{vous} et passe au \emph{tu} au milieu d'une ligne,
à « J'ai à te dire », pour n'y plus revenir — sauf dans la formule finale, qui
est de forme et reprend le \emph{votre}.\footnote{Le fait est noté dans la
transcription à l'endroit même où il se produit. On le relève ici parce qu'il
marque le moment où la lettre cesse d'être une lettre de nouvelles pour
devenir une réponse personnelle : le \emph{vous} couvre l'excuse du silence,
le \emph{tu} commence avec le reproche.}

\section*{Le billet du père (vue 15)}

\pagerange{15}{15}
Entre la fin de la lettre et sa formule finale, dans le blanc laissé au milieu
de la page, une seconde main a écrit treize lignes. Elle est plus petite,
beaucoup plus rapide, à longs traits horizontaux, et par endroits pâlie ;
plusieurs de ses mots n'ont pas été lus. Elle appelle la destinataire « ma
bonne fille » et signe « votre bon père » — c'est donc le père, Étienne
Pascal, écrivant à sa fille sur la lettre de son fils.\footnote{Le feuillet ne
donne, là encore, que le nom « Pascal », en grand, au milieu de la page, avec
le cachet rond de la bibliothèque qui le chevauche. L'identification tient à
la formule « V\textsuperscript{re} bon Pere » et à ce que la lettre principale
parle de « mon Pere » comme présent auprès de son auteur.}

Ma bonne fille mignonne, si vous savez que je ne prends pas le temps
d'écrire,\footnote{La fin de la première ligne du billet, quatre ou cinq mots
entre « Si Vous » et « d'escrire », n'a pas été lue ; la phrase donnée ici
n'est pas une restitution mais un raccord, et le début du billet reste
inconnu.} c'est que je n'ai aucun loisir : car je n'ai jamais été dans
l'embarras à la dîme de ce que j'y suis à présent, et je ne saurais l'être
davantage à moins d'y avoir trop.\footnote{Le feuillet a « a la
des\textsuperscript{me} », suivi d'un mot d'une dizaine de jambages non lu.
Cette lecture propose « à la dîme », c'est-à-dire au dixième, parce que la
phrase demande une fraction et que la construction « n'avoir jamais été dans
l'embarras à la dîme de ce que j'y suis » se comprend alors immédiatement : je
n'ai jamais été dix fois moins occupé qu'aujourd'hui, ce qui est la forme
ordinaire de l'hyperbole. La proposition est de cette lecture ; le mot suivant
reste illisible et pourrait changer le sens.} Il y a quatre mois que je ne me
suis pas couché six fois avant deux heures après
minuit.\footnote{« Quatre » est écrit au-dessus de « trois » biffé : le père a
corrigé son compte vers le haut. Sur quatre mois, soit environ cent vingt
nuits, moins de six nuits font moins d'une nuit sur vingt. Le calcul est de
cette lecture ; le feuillet ne donne que les deux nombres.}

Je vous envoie une lettre sur ce sujet de la lettre dernière, touchant les
bruits du mariage de monsieur Despeaux ; mais je n'ai jamais su que ce fût de
l'archevêque.\footnote{Trois lacunes dans cette phrase. Entre « envoy » et
« Une l\textsuperscript{re} », deux ou trois mots non lus. Le nom de
l'expéditeur de cette lettre incluse, après « de », compte huit ou neuf
lettres, commence par une majuscule qui se lit B ou R et paraît finir en
\emph{-bre} : la transcription ne le donne pas, et cette lecture non plus.
« Despeaux » et « que ce fût » sont donnés avec doute.}

Pour nouvelles : la fille de monsieur de Paris, madame du Comptes, mariée à
monsieur de Montiffer ; aussi monsieur du Comptes est décédé ; comme aussi la
fille de Belair, mariée au petit Lambert.\footnote{Les noms propres de ce
paragraphe sont écrits très vite et sans reprise, et la transcription les
marque tous douteux ; ils sont repris ici tels quels et ne doivent pas être
tenus pour lus. Le feuillet a « est decedée » au féminin après
« M\textsuperscript{r} du Comptes » ; l'accord est rétabli au masculin
ci-dessus, mais rien n'assure que le père n'ait pas voulu écrire que c'est la
dame qui est morte, auquel cas la phrase se contredirait avec le mariage qui
précède.} Je suis toujours votre bon père et très affectionné ami,

Pascal.

Votre petit a couché cette nuit ; il se porte, Dieu grâce, très
bien.\footnote{Trois lignes serrées dans la marge inférieure gauche, de la
même main, en regard de la formule finale ; un mot bref à la fin de la
première ligne n'a pas été lu, et « ceste » est donné avec doute. L'enfant
n'est pas nommé : c'est un petit-fils ou une petite-fille du signataire, et
donc un enfant de la destinataire, qui se trouve auprès de son grand-père et
non auprès de sa mère. Le feuillet n'en dit pas plus.}

Le billet est de la même veine que la lettre : il mesure. Le père ne dit pas
qu'il est débordé, il dit de combien — d'un facteur dix pour la charge, de six
nuits sur cent vingt pour le sommeil. Et il finit, comme la lettre de son
fils, sur ce qui manque : des nouvelles d'une malade, la santé d'un enfant.

\section*{L'adresse (vue 16)}

\pagerange{16}{16}
Au dos du second feuillet, écrit de la main de la lettre au milieu de la page,
le feuillet ayant été plié :

À Mademoiselle, mademoiselle Périer, la conseillère, à
Clermont.\footnote{Sur quatre lignes, les deuxième et troisième en retrait.
« La conseillère » désigne la femme du conseiller, selon l'usage qui donne à
l'épouse le féminin de la charge du mari ; « Mademoiselle » est le titre des
femmes mariées de la bourgeoisie et de la petite noblesse, et ne dit rien du
célibat. Aucune marque postale n'est lisible.}

\section*{La quittance de 1674 (vues 9 et 10)}

\pagerange{9}{10}
Le feuillet à l'italienne des vues 9 et 10 n'a rien à voir avec la lettre et
n'est pas mentionné dans la notice. C'est un formulaire gravé, rempli d'une
main de notaire : une quittance pour les rentes de l'Hôtel de Ville de
Paris.\footnote{Le titre gravé est « Quittance pour les Rentes de l'Hostel de
Ville de Paris », coupé en son milieu par un cartouche ovale portant en cercle
« GENERALITE DE PARIS » et une fleur de lys. La mention « Deux sols », en
tête, est le droit du formulaire. En marge, en haut et souligné : « 1674.
derniers six moys ».}

\subsection*{Le texte}

Par-devant les notaires royaux établis en la ville de Clermont fut présente
dame Gilberte Pascal, veuve et héritière testamentaire de défunt
M\textsuperscript{e} Florin Périer, vivant conseiller du roi en sa Cour des
aides de Clermont-Ferrand, laquelle de son bon gré a confessé avoir reçu
la somme de cent trente-trois livres quatre sols, ordonnée être laissée en
fonds dans l'état du roi de la ferme générale des aides et entrées pour les
six derniers mois de l'année présente mille six cent soixante-quatorze, à
cause de huit cent quatre-vingt-huit livres de rente constituée sur les
tailles par contrat du second janvier mille six cent trente-cinq, et à présent
assignée sur lesdites aides et entrées suivant la déclaration de Sa Majesté du
treizième janvier mille six cent soixante-cinq.\footnote{Le nom du payeur,
après « avoir receu de », est resté en blanc sur le formulaire, et la ligne
s'arrête là ; cette lecture enchaîne donc directement sur la somme, ce que la
phrase du formulaire ne fait pas. Une \emph{rente constituée} est une rente
perpétuelle créée par contrat contre un capital versé ; elle est ici
\emph{assignée} sur un revenu fiscal, c'est-à-dire que le paiement en est
gagé sur le produit des aides et entrées, la déclaration de 1665 ayant déplacé
ce gage. L'\emph{état du roi} est le budget arrêté d'où le paiement doit
sortir. La \emph{Cour des aides} était la juridiction souveraine du contentieux
fiscal.} De laquelle somme ladite dame quitte ledit sieur payeur et tous
autres. Fait et passé audit Clermont, étude de l'un desdits notaires, le
vingt-septième jour de juin mille six cent soixante-quatorze, après
midi.\footnote{Les formules d'usage abrégées du feuillet — « Ainsy \&c. A
peine \&c. Oblige \&c. Renonce \&c. Soumet \&c. » — sont omises ici ; elles
renvoient aux clauses ordinaires de l'acte notarié. Un long trait de
remplissage suit « vingt septiesme » pour interdire l'ajout d'un mot.} Ladite
dame confessante a signé avec lesdits notaires. J'approuve le mot de
six.\footnote{« Six », dans « six derniers moys », est ajouté au-dessus de la
ligne ; la mention finale l'approuve, selon la règle qui veut qu'un mot
interligné soit couvert par une approbation signée, faute de quoi il est
réputé non écrit. La fin de cette approbation, en petite cursive rapide, n'a
pas été lue.}

Signé : Camus, notaire royal à Clermont ; G. Pascal ; Guillotteau, notaire
royal.\footnote{Les deux notaires signent avec un grand paraphe. La signature
« G. Pascal. » est autographe de Gilberte Pascal et se trouve au milieu de la
ligne des signatures ; le cachet rond « BIBLIOTHÈQUE NATIONALE MSS » est
apposé juste au-dessous. Au bas du texte, une ligne de dos — « Quittance de
cent trente-trois livres quatre sols » — en partie recouverte par les
paraphes.}

Au verso, écrit en haut à droite, un certificat d'une semaine postérieur :
« Je soussigné Antoine Ferret, bourgeois de Paris, certifie à tous qu'il
appartiendra que la quittance de l'autre part est véritable. À Paris, ce
quatrième juillet mille six cent soixante-quatorze. »\footnote{Le nom est
donné d'après la signature qui suit, et la transcription le marque douteux :
l'initiale pourrait aussi bien se lire L. Le millésime est écrit en abrégé,
deux ou trois lettres surmontées d'un c ; la lecture « mil vj\textsuperscript{c} »
est offerte par le sens, « soixante-quatorze » étant net. La pièce est passée
à Clermont et certifiée à Paris : la quittance doit remonter au payeur
parisien, et le certificat d'un bourgeois de Paris supplée à ce que les
notaires d'Auvergne n'y sont pas connus.}

\subsection*{Le compte de la rente}

La pièce donne deux nombres et ne les rapproche jamais : une rente de huit
cent quatre-vingt-huit livres par an, et cent trente-trois livres quatre sols
reçues pour un semestre. Le semestre dû serait de quatre cent quarante-quatre
livres. Or, la livre valant vingt sols, cent trente-trois livres quatre sols
font $133 + \tfrac{4}{20} = 133{,}2$ livres, et
\[
\frac{133{,}2}{444} = \frac{3}{10}.
\]
Le paiement est donc exactement trois dixièmes de ce que la rente porte, ou,
sur l'année, deux cent soixante-six livres huit sols au lieu de huit cent
quatre-vingt-huit.\footnote{La division est de cette lecture ; le feuillet ne
la fait pas et n'énonce aucun taux. Le résultat tombe rond, ce qui indique un
abattement fixé en fraction et non un calcul de prorata : la monarchie a
plusieurs fois \emph{retranché} les rentes sur l'Hôtel de Ville, c'est-à-dire
réduit d'autorité ce qu'elle en servait, et la pièce est le reçu de ce qui
reste. Rien dans le feuillet ne nomme cet abattement ni n'en donne la date ;
la lecture s'en tient au rapport des deux nombres.} On remarquera que le
compte tombe juste au sol près : le cinquième d'une livre fait exactement
quatre sols dans une monnaie de compte où la livre en vaut vingt. Le même
diviseur cinq aurait donné un reste dans une monnaie décimale.\footnote{La
monnaie de compte est la livre tournois, valant vingt sols de douze deniers,
soit deux cent quarante deniers. Un système de ce genre divise exactement par
2, 3, 4, 5, 6, 10, 12, 15 et 20, ce que le système décimal ne fait ni pour 3
ni pour 12 : c'est la raison pour laquelle les monnaies d'Ancien Régime
étaient bâties ainsi, et c'est aussi pourquoi un abattement de trois dixièmes
s'écrit ici sans fraction résiduelle.}

Ce que la pièce ne permet pas de calculer, il faut le dire : le capital de la
rente n'y est pas, donc le taux — le \emph{denier} auquel elle fut constituée
en 1635 — ne se retrouve pas. On sait seulement qu'elle a été créée par
contrat du 2 janvier 1635, d'abord gagée sur les tailles, puis reportée sur
les aides et entrées en 1665, et qu'en 1674 elle appartenait à la succession
de Florin Périer.

\subsection*{Le chiffre en vingtaines}

Au milieu du verso, à gauche de la signature, le feuillet porte un nombre
écrit à l'ancienne notation par vingtaines, dont la transcription ne lit que
la fin : un caractère illisible, puis \textsuperscript{xx}, puis xiij. Cette
notation écrit un nombre comme un multiple de vingt plus un reste, le
multiplicateur étant placé en exposant sur le xx. Comme $133 = 6 \times 20 +
13$, le nombre s'écrit vj\textsuperscript{xx} xiij, et le caractère perdu est
un vj.\footnote{La restitution est de cette lecture et repose sur
l'arithmétique seule : la transcription note que le nombre s'accorderait avec
les cent trente-trois livres de la quittance, mais se refuse à faire le
rapprochement, la pièce ne le disant pas. Le rapprochement est fait ici parce
qu'aucun autre nombre du feuillet ne finit par treize et que les deux chiffres
lisibles, \textsuperscript{xx} et xiij, ne laissent qu'un multiplicateur
possible une fois la somme connue. C'est un argument, non une lecture : si
l'on tenait le nombre pour indépendant de la somme, le multiplicateur
resterait inconnu.} C'est le montant de la quittance, répété en chiffres à
côté des signatures, comme les actes le font pour interdire la surcharge.

Cette notation est la même vingtaine qui structure la monnaie du feuillet, et
la même que le français a gardée dans « quatre-vingts » : le nombre huit cent
quatre-vingt-huit de la rente en est l'exemple, écrit en toutes lettres deux
lignes plus haut.\footnote{Au bas du verso, tête-bêche, quatre lignes très
pâles — un millésime 1674, le mot « escheue », un mois et des chiffres — n'ont
pas été lues ; elles sont probablement la mention d'échéance et pourraient
porter d'autres nombres.}

\section*{Ce que ce volume n'est pas}

\pagerange{9}{16}
Il faut finir par ce qui manque, parce que la cote invite à l'erreur. Ces
feuillets ne contiennent aucune mathématique de Blaise Pascal : ni calcul, ni
figure, ni raisonnement autre que le petit compte des lettres perdues. La
lettre de 1643 est de l'année même où il travaillait à la machine
arithmétique, et le motif de cette machine était précisément le travail dont
la lettre parle — les comptes de la levée des tailles en Normandie, la tournée
des commissaires, les rôles à établir. La lettre n'en dit pas un mot, et cette
lecture ne prêtera pas au feuillet ce qu'il ne porte pas.\footnote{La
chronologie de la machine n'est pas sur le feuillet et n'est pas discutée ici.
On note seulement que les deux choses se rencontrent dans les mêmes semaines
et dans le même bureau, et qu'il serait faux d'en tirer que la lettre y
fasse allusion : elle ne parle que du courrier et d'une malade.}

La quittance de 1674, elle, contient bien un calcul — mais ce n'est pas celui
d'un mathématicien. C'est l'arithmétique d'un bureau : une fraction imposée
d'en haut, appliquée à un principal, et le résultat porté deux fois, en
lettres dans le corps de l'acte et en chiffres à côté des signatures. Ce qui
la rend intéressante dans ce volume est qu'elle est signée de la main de la
femme à qui la lettre était adressée, et qu'entre les deux pièces il s'est
passé une vie entière.

\end{document}
